\section{Introduction}
\label{sec:intro}

An office accumulates idle computation. Desktops sit at a few percent
utilisation through the working day and at zero overnight; small machines bought
for one task spend the rest of their lives idle. The aggregate is substantial,
and none of it is reachable by anyone who needs compute.

Batch systems solve exactly this, and HTCondor was built for the case: harvesting
cycles from machines whose owners have first claim on them. The engine is not the
problem. The problem is the interaction model that conventionally arrives with
it.

\begin{keypoint}
The requirement that shaped every decision here is that users must never log
into a submission host. Not because logging in is difficult, but because the
moment it is required the batch system becomes a separate place people have to
visit, with its own home directory, its own environment and its own files to
copy back and forth, rather than an extension of the machine they already work
on.
\end{keypoint}

That friction is what keeps shared clusters underused. A colleague with a script
that would finish in ten minutes across a pool will run it for six hours locally
rather than learn another machine's filesystem layout.

\subsection{Why HTCondor}

Desktop grids and volunteer computing systems are mature, and the field is
large. The distinction that matters for this deployment is where generality
sits. Volunteer computing platforms are typically built for one application or
one group and generalised afterwards; participants run a client that pulls work
units, and the scheduling model is the platform's own.

HTCondor starts general. ClassAd matchmaking makes machine policy a first-class,
user-editable expression rather than a platform behaviour. Owner priority is a
founding feature rather than an addition. And it is indifferent to UID
consistency across hosts, which matters here because the same login has a
different UID on different machines.

The alternatives fail on specifics rather than in principle. Slurm assumes a
consistent user namespace across the cluster, which this fleet does not have.
Kubernetes solves container orchestration for services, and its eviction model
is driven by capacity pressure rather than by the arrival of a human at a
keyboard. A dedicated cluster is the right answer when the hardware exists; the
premise here is that the machines are already bought, already deployed, and
already doing something else most of the time.

The consequence is that this work is not the construction of a system. It is the
discovery of which of an existing system's controls do what they claim, and the
policy written on top of them.

\subsection{Contributions}

\begin{enumerate}
  \item \textbf{Measurement-based admission control} for an opportunistic pool,
        with every gating term drawn from \file{/proc} rather than from what
        jobs claimed they would use, including Linux pressure stall information
        as the signal that an owner is actually suffering
        (\S\ref{sec:admission}).

  \item \textbf{A demonstration that the reservation ledger, not the policy, is
        the binding gate}, and that it can be replaced with measured headroom.
        The matchmaker consults the advertised \kn{Memory} attribute before any
        policy expression runs, so a machine whose jobs over-request refuses
        work while idle. Measured before and after (\S\ref{sec:ledger}).

  \item \textbf{Naming and addressing for execute nodes that move between
        unroutable subnets}, where the worker holds no address, no network name
        and no per-machine configuration, and derives at startup which of the
        entry node's two advertised addresses it can reach
        (\S\ref{sec:mobility}).

  \item \textbf{A catalogue of eleven silent failures} sharing one shape, with a
        structural account of why declarative policy systems produce that shape
        and the verification discipline that follows (\S\ref{sec:silent}).
\end{enumerate}

\subsection{What was built}

A four-machine pool spanning two subnets with no route between them: an entry
node carrying the collector, negotiator, schedd and a limited startd; a
sixteen-thread worker; and two four-core desktops in daily use by their owners,
one of which is also the submit client for its owner's own jobs. Jobs submitted
from a user's own desktop distribute across all four machines and return their
output. All measurements in this note come from that pool.
